\section{Consistency equivalence, reproducibility, and limits}
\label{sec:verification}

\begin{theorem}[The requested equivalence]\label{thm:consistency}
Let the fixed four integers be the tuple in
Section~\ref{sec:coefficients}. ZFC proves
\[
\begin{split}
\neg\exists\mathbf x\in\Z_{>0}^{12}\,
 P_{12}(\mathbf x)=0
&\ \longleftrightarrow\ \Con(\mathrm{ZFC}),\\
\neg\exists\mathbf z\in\Z^{24}\,
 P_{24}(\mathbf z)=0
&\ \longleftrightarrow\ \Con(\mathrm{ZFC}).
\end{split}
\tag{8.1}
\]
The same equivalences hold for the shifted three-equation
systems (5.14), (6.14), in their stated domains, and for the
integer sum of squares (6.15).
\end{theorem}
\begin{proof}
Theorems~\ref{thm:positive} and \ref{thm:integer}, with the
coefficient adapter, give finite polynomial solutions if and
only if there is a finite blank halting computation of the
literal machine. The exact machine calibration identifies those
computations with finite ZFC refutations.
All reductions use finite coding, integer arithmetic, Pell
recurrences, finite algebraic extensions, and the indicated
classical number-theoretic facts; their proofs can be carried
out in ZFC.
Thus ZFC proves each existential statement equivalent to its
syntactic refutability statement. Negating gives (8.1).
The system shifts preserve existence in both directions,
and over the integers a sum of squares vanishes exactly when
every summand vanishes.
\end{proof}

This theorem does not decide which side holds.
In particular, under the usual consistency assumption, proving
one of these polynomials has no solutions within ZFC would
contradict the second incompleteness theorem.
The construction identifies the arithmetic statement whose
unprovability is at issue.

\subsection{What is literal and what is inherited}

The archive is intended as a reviewable repository contribution.
It contains:
\begin{itemize}
\item this \LaTeX{} source and its PDF, with the new proofs and
  exact nonvanishing degree certificates;
\item seven literal JSON arithmetic circuits: three for the
  positive construction and four for the integer construction;
\item Python generators and finite arithmetic diagnostics,
  plus an independent circuit/liveness/degree-bound checker;
\item the complete coefficient recipe, transition table, ordered
  Wang tile list, four compressed binary integers, and independent
  byte-for-byte reconstruction checks;
\item independent written reviews and the source copies needed
  to audit the reused norm, signed-domain, compiler, and endpoint
  interfaces;
\item source provenance, licenses, an artifact checksum manifest,
  and a root reproduction command.
\end{itemize}

The circuit files enumerate the supplied variables and the four
fixed ports separately, define every intermediate using a binary
operation, and identify their output equation or comparisons.
They can be expanded mechanically to ordinary polynomials,
although doing so loses useful sharing and greatly enlarges the
presentation. The degrees in Table~\ref{tab:new} are the degrees
of those expanded polynomials, not circuit-depth estimates.

The newly supplied proofs establish the two root reductions and
the coefficient adapter. The inherited native compiler,
Pell-index lemmas, Wang language theorem, and machine/refutation
calibration remain dependencies. Their pinned sources are
identified and relevant copies are included.
We do not claim to have independently reconstructed every
earlier compiler proof or rerun the complete formal endpoint.

\subsection{Reproduction and the scope of the checks}

From the archive root, with a standard Python 3 installation,
run
\begin{quote}\ttfamily
python3 reproduce.py
\end{quote}
No network, computer algebra package, or proof assistant is
required for the finite checks.
To rebuild the manuscript, use a standard \TeX{} installation:
\begin{quote}\ttfamily
pdflatex -interaction=nonstopmode -halt-on-error report.tex\\
pdflatex -interaction=nonstopmode -halt-on-error report.tex
\end{quote}

The arithmetic checks compare the compact norm evaluator with
an independent eight-dimensional conjugate-product evaluation
on signed assignments; exercise the inverse reconstruction
including negative strong roots; verify transport shifts;
audit live gates and port use; and check nonvanishing leading
forms modulo two primes.
The integer review separately derives the signed bounds,
the two boundary exclusions, and the exact weighted leading form.
The coefficient checks cover \(69{,}230\) basis convolution
identities, all-rotation support facts, carry/slack inequalities,
population counts, the five-adic unit, and the four exact payloads.
The inherited language checker tests \(425\) explicitly halting
control helices comprising \(33{,}800\) local tiles.

These tests detect transcription errors and verify the named
finite properties. Finite fixtures do not establish a universal
Diophantine equivalence. That conclusion depends on the proofs,
including the inherited compiler, rather than on a search for
zeros or a successful sample of computations.
No new Lean/Rocq proof object for (8.1) is supplied.

\subsection{Which presentation to use}

For twelve positive unknowns, (5.14) is the recommended system:
three equations, \(110\) operations, maximum degree \(1072\).
For a single equation with that domain, use (5.2).
For ordinary unrestricted integer unknowns, use (6.14) if a
system is acceptable and (6.4) if one equation is required.
The sum of squares is useful when its especially direct
zero-set argument outweighs its higher degree.

The contribution improves the inspected moderate-degree
twelve-positive and twenty-six-integer families.
It does not dominate the retained \(67\)-operation positive
polynomial, the \(59\)-operation larger-variable system,
the degree-\(172\) twenty-seven-integer product, or the
nine-natural-variable Jones specialization.
There is no proved global optimum among all Diophantine
representations of this consistency statement.

A further variable-count comparison is timely.
Sun's September 2026 preprint \cite{sun10}, Theorem 1.1,
states that every recursively enumerable subset of \(\N\)
has a representation using ten integer existential unknowns.
Its theorem gives a generic route to a ten-integer consistency
instance. We have not instantiated or independently audited
that compiler here, so this report makes no claim to a global
integer-variable record.

\subsection{Suggested next work}

The best immediate verification target is to formalize the
new signed separation and saturation composition as a small
extension of the existing core interface. The positive proof's
new obligation is similarly localized to (5.5)--(5.8).
For further size reduction, the most promising open comparison
is whether a weaker integer-valued or nonmonic inverse
construction can remove another guarded coordinate without
destroying soundness at zero or negative scales.
Any such proposal must recheck coherent root recovery,
the \(C=0\) boundary, exact leading cancellation, and simultaneous
canonical completeness. The present work stops at the improvements
proved and reproduced in this archive.
